\documentclass[../Main.tex]{subfiles}
\begin{document}
  \chapter{2015-2016学年微积分（一）（下）期末考试}

  \section{单项选择题(每小题 3 分，共 18 分)}
  \begin{enumerate}
    \item 曲面$z=4-x^{2}-y^{2}$上点$P$处的切平面平行于平面$2x+2y+z=1$，则点$P$的坐标是（\qquad）.

      \begin{tasks}[label=\textbf{\Alph*.}, label-width=1.5em, column-sep=1em](2)
        \task $(1,-1,2)$
        \task $(-1,1,2)$
        \task $(1,1,2)$
        \task $(-1,1,1)$
      \end{tasks}
      \vspace{1em}

    \item 考虑二元函数$f(x,y)$在点$(x_{0},y_{0})$处的以下 4 条性质：

      （1）两个偏导数存在；（2）两个偏导函数连续；（3）可微；（4）沿任意方向的方向导数存在。

      若用“$P\Rightarrow Q$”表示可由性质$P$推出性质$Q$，则有（\qquad）.

      \begin{tasks}[label=\textbf{\Alph*.}, label-width=1.5em, column-sep=1em](2)
        \task $(2)\Rightarrow(3)\Rightarrow(1)$
        \task $(3)\Rightarrow(2)\Rightarrow(1)$
        \task $(3)\Rightarrow(4)\Rightarrow(1)$
        \task $(3)\Rightarrow(1)\Rightarrow(4)$
      \end{tasks}
      \vspace{1em}

    \item 将二次积分
      \[
        I=\int_{0}^{\frac{\pi}{4}}\,\mathrm{d}\theta\int_{0}^{1}f(r\cos\theta,r\sin\theta)r\,\mathrm{d}r
      \]
      化为直角坐标系下的二次积分，正确的结果是（\qquad）.

      \begin{tasks}[label=\textbf{\Alph*.}, label-width=1.5em, column-sep=1em](1)
        \task $\displaystyle I=\int_{0}^{1}\,\mathrm{d}y\int_{y}^{\sqrt{1-y^{2}}}f(x,y)\,\mathrm{d}x$
        \task $\displaystyle I=\int_{0}^{1}\,\mathrm{d}x\int_{\sqrt{1-x^{2}}}^{x}f(x,y)\,\mathrm{d}y$
        \task $\displaystyle I=\int_{0}^{\frac{\sqrt2}{2}}\,\mathrm{d}y\int_{y}^{\sqrt{1-y^{2}}}f(x,y)\,\mathrm{d}x$
        \task $\displaystyle I=\int_{0}^{\frac{\sqrt2}{2}}\,\mathrm{d}x\int_{\sqrt{1-x^{2}}}^{x}f(x,y)\,\mathrm{d}y$
      \end{tasks}
      \vspace{1em}

    \item 设$S$是圆柱面$x^{2}+y^{2}=R^{2}(R>0)$被平面$z=0,z=h$截下的部分，则
      \[
        \iint\limits_{S}(x^{2}+y^{2})\,\mathrm{d}S=(\qquad).
      \]

      \begin{tasks}[label=\textbf{\Alph*.}, label-width=1.5em, column-sep=1em](2)
        \task $-2\pi hR^{3}$
        \task $2\pi hR^{3}$
        \task $\pi hR^{3}$
        \task $\dfrac{\pi}{2}R^{4}h$
      \end{tasks}
      \vspace{1em}

    \item 下列命题中正确的是（\qquad）.

      \begin{tasks}[label=\textbf{\Alph*.}, label-width=1.5em, column-sep=1em](1)
        \task 若$\lim\limits_{n\to\infty}a_{n}=0$，则$\sum\limits_{n=1}^{\infty}a_{n}$收敛
        \task 若正项级数$\sum\limits_{n=1}^{\infty}a_{n}$收敛，且$l=\lim\limits_{n\to\infty}\dfrac{a_{n+1}}{a_{n}}$存在，则$l<1$
        \task 若$\lim\limits_{n\to\infty}\dfrac{a_{n}}{b_{n}}=1$且$\sum\limits_{n=1}^{\infty}b_{n}$收敛，则$\sum\limits_{n=1}^{\infty}a_{n}$收敛
        \task 若$\sum\limits_{n=1}^{\infty}a_{n}$绝对收敛，则$\sum\limits_{n=1}^{\infty}a_{2n}$绝对收敛
      \end{tasks}
      \vspace{1em}

    \item 设$f(x)=x+1(0\leq x\leq \pi)$，而
      \[
        S(x)=\frac{a_{0}}{2}+\sum_{n=1}^{\infty}a_{n}\cos nx,\qquad -\infty<x<+\infty,
      \]
      其中
      \[
        a_{n}=\frac{2}{\pi}\int_{0}^{\pi}f(x)\cos nx\,\mathrm{d}x,
      \]
      则$S\left(-\frac{\pi}{2}\right)=(\qquad)$.

      \begin{tasks}[label=\textbf{\Alph*.}, label-width=1.5em, column-sep=1em](2)
        \task $-\dfrac{\pi}{2}+1$
        \task $\dfrac{\pi}{2}+1$
        \task $-\dfrac{\pi}{2}-1$
        \task $\dfrac{\pi}{2}-1$
      \end{tasks}
      \vspace{1em}
  \end{enumerate}

  \section{填空题(每小题 4 分，共 16 分)}
  \begin{enumerate}
    \setcounter{enumi}{6}
    \item 以曲线
      \[
        \begin{cases}
          2x^{2}+y^{2}+z^{2}=16, \\
          x^{2}+y^{2}-z^{2}=0
        \end{cases}
      \]
      为准线，母线平行于$x$轴的柱面方程是\underline{\hspace{4cm}}.

    \item 设$F=\{y,z,xy\}$，$G=\{1,0,x\}$，则$\operatorname{rot}(F\times G)=$\underline{\hspace{3cm}}.

    \item 设$D:x^{2}+y^{2}\leq a^{2}(a>0)$，若
      \[
        \iint\limits_{D}\sqrt{a^{2}-x^{2}-y^{2}}\,\mathrm{d}x\,\mathrm{d}y=\pi,
      \]
      则$a=$\underline{\hspace{3cm}}.

    \item 设$L$为从$A(1,0)$到$B(2,1)$的曲线弧，则
      \[
        \int_{L}(2xy-y^{4}+3)\,\mathrm{d}x+(x^{2}-4xy^{3})\,\mathrm{d}y=
      \]
      \underline{\hspace{3cm}}.
  \end{enumerate}

  \section{基本计算题(每小题 7 分，共 42 分)}
  \begin{enumerate}
    \setcounter{enumi}{10}
    \item 求过点$P(1,2,3)$且与$z$轴相交、与平面$x+y+z-2=0$平行的直线方程.
      \vspace{11em}

    \item 设$z=f\left(\sin x,\dfrac{x}{y}\right)$，其中$f$具有二阶连续导数，求$\dfrac{\partial^{2}z}{\partial x\partial y}$.
      \vspace{12em}

    \item 设函数$y=y(x),z=z(x)$由方程$z=xf(x+y)$以及$F(x,y,z)=0$所确定，其中$f$具有一阶连续导数，$F$具有一阶连续偏导数，求$\dfrac{\mathrm{d}z}{\mathrm{d}x}$.
      \vspace{13em}

    \item 求三重积分
      \[
        I=\iiint\limits_{V}(x^{2}+yz+z^{2})\,\mathrm{d}v,
      \]
      其中$V:x^{2}+y^{2}+z^{2}\leq1$.
      \vspace{12em}

    \item 求曲线积分
      \[
        I=\oint\limits_{L}(x^{2}-y^{2})\,\mathrm{d}x+(x^{2}+y^{2})\,\mathrm{d}y,
      \]
      其中$L$是圆周$x^{2}+y^{2}=x+y+1$，取正向.
      \vspace{12em}

    \item 求幂级数$\displaystyle\sum_{n=1}^{\infty}\dfrac{(x-3)^{n}}{n\cdot3^{n}}$的收敛域与和函数$S(x)$.
      \vspace{13em}
  \end{enumerate}

  \section{应用题(每小题 7 分，共 14 分)}
  \begin{enumerate}
    \setcounter{enumi}{16}
    \item 在椭球面$x^{2}+2y^{2}+2z^{2}=1$上求一点，使函数$f(x,y,z)=x^{2}+y^{2}+z^{2}$在该点处沿方向$\boldsymbol{n}=\{0,-1,1\}$的方向导数最大.
      \vspace{12em}

    \item 设速度场$\boldsymbol{v}=\{2x^{3},2y^{3},3(z^{2}-1)\}$，$S$是曲面$z=1-x^{2}-y^{2}(z\geq0)$的上侧，求单位时间内流过曲面$S$（上侧）的流量（或通量）
      \[
        \Phi=\iint\limits_{S}\boldsymbol{v}\cdot\boldsymbol{n}\,\mathrm{d}S.
      \]
      \vspace{12em}
  \end{enumerate}

  \section{综合题(每小题 5 分，共 10 分)}
  \begin{enumerate}
    \setcounter{enumi}{18}
    \item 证明莱布尼兹判别法：若数项级数$\displaystyle\sum_{n=1}^{\infty}(-1)^{n-1}a_{n}(a_{n}>0)$满足（1）$\lim\limits_{n\to\infty}a_{n}=0$；（2）$\{a_{n}\}$单调减少，则级数$\displaystyle\sum_{n=1}^{\infty}(-1)^{n-1}a_{n}$收敛.
      \vspace{15em}

    \item 设函数$f(x,y)$具有二阶连续偏导数，且$f(1,y)=0,f(x,1)=0,\iint\limits_{D}f(x,y)\,\mathrm{d}\sigma=a$，其中$D:0\leq x\leq1,0\leq y\leq1$。计算积分
      \[
        I=\iint\limits_{D}xyf_{xy}(x,y)\,\mathrm{d}\sigma.
      \]
      \vspace{15em}
  \end{enumerate}

  \chapter{2015-2016学年微积分（一）（下）期末考试参考答案}
  \section{单项选择题(每小题 3 分，共 18 分)}
  \begin{enumerate}
    \item \textbf{Solution}. C.
    \item \textbf{Solution}. A.
    \item \textbf{Solution}. C.
    \item \textbf{Solution}. B.
    \item \textbf{Solution}. D.
    \item \textbf{Solution}. B.
  \end{enumerate}

  \section{填空题(每小题 4 分，共 16 分)}
  \begin{enumerate}
    \setcounter{enumi}{6}
    \item \textbf{Solution}. $3z^{2}-y^{2}=16$.
    \item \textbf{Solution}. $\{0,x,0\}$.
    \item \textbf{Solution}. $\sqrt[3]{\dfrac{3}{2}}$.
    \item \textbf{Solution}. $5$.
  \end{enumerate}

  \section{基本计算题(每小题 7 分，共 42 分)}
  \begin{enumerate}
    \setcounter{enumi}{10}
    \item \textbf{Solution}. 设所求直线的方向向量为$\boldsymbol{s}$，并与$z$轴相交于点$Q(0,0,c)$，则
      \[
        \boldsymbol{s}=\overrightarrow{PQ}=\{-1,-2,c-3\}.
      \]
      由题设知$\boldsymbol{s}\perp\{1,1,1\}$，故$c=6$.
      所求直线方程为
      \[
        \frac{x-1}{1}=\frac{y-2}{2}=\frac{z-3}{-3}.
      \]

    \item \textbf{Solution}. 记$f_{1},f_{2},f_{12},f_{22}$分别表示$f$对其变量的一阶、二阶偏导数，则
      \[
        z_{x}=\cos x\,f_{1}+\frac{1}{y}f_{2}.
      \]
      因而
      \[
        z_{xy}=-\frac{x\cos x}{y^{2}}f_{12}-\frac{1}{y^{2}}f_{2}-\frac{x}{y^{3}}f_{22}.
      \]

    \item \textbf{Solution}. 对方程组
      \[
        \begin{cases}
          z=xf(x+y), \\
          F(x,y,z)=0
        \end{cases}
      \]
      两边对$x$求导，得
      \[
        \begin{cases}
          \dfrac{\mathrm{d}z}{\mathrm{d}x}=f+x\left(1+\dfrac{\mathrm{d}y}{\mathrm{d}x}\right)f', \\[8pt]
          F_{1}+\dfrac{\mathrm{d}y}{\mathrm{d}x}F_{2}+\dfrac{\mathrm{d}z}{\mathrm{d}x}F_{3}=0.
        \end{cases}
      \]
      消去$\dfrac{\mathrm{d}y}{\mathrm{d}x}$可得
      \[
        \frac{\mathrm{d}z}{\mathrm{d}x}
        =\frac{(f+xf')F_{2}-xf'F_{1}}{F_{2}+xf'F_{3}}.
      \]

    \item \textbf{Solution}. 由球域对称性，
      \[
        \iiint\limits_{V}yz\,\mathrm{d}v=0,\qquad
        \iiint\limits_{V}(x^{2}+yz+z^{2})\,\mathrm{d}v
        =\frac{2}{3}\iiint\limits_{V}(x^{2}+y^{2}+z^{2})\,\mathrm{d}v.
      \]
      改用球坐标，
      \[
        I=\frac{2}{3}\int_{0}^{2\pi}\,\mathrm{d}\theta\int_{0}^{\pi}\sin\varphi\,\mathrm{d}\varphi\int_{0}^{1}\rho^{4}\,\mathrm{d}\rho
        =\frac{8\pi}{15}.
      \]

    \item \textbf{Solution}. 由格林公式，
      \[
        I=\iint\limits_{D}\left[\frac{\partial(x^{2}+y^{2})}{\partial x}-\frac{\partial(x^{2}-y^{2})}{\partial y}\right]\mathrm{d}x\,\mathrm{d}y
        =\iint\limits_{D}2(x+y)\,\mathrm{d}x\,\mathrm{d}y.
      \]
      圆周$x^{2}+y^{2}=x+y+1$所围区域$D$的圆心为$\left(\dfrac12,\dfrac12\right)$，半径为$\sqrt{\dfrac32}$，
      故
      \[
        I=2(\bar{x}+\bar{y})\sigma=2\cdot1\cdot\frac{3\pi}{2}=3\pi.
      \]

    \item \textbf{Solution}. 令$t=x-3$，则原级数化为
      \[
        \sum_{n=1}^{\infty}\frac{t^{n}}{n3^{n}}.
      \]
      收敛半径为$3$，且当$t=3$时发散，当$t=-3$时收敛，
      故原级数收敛域为
      \[
        [0,6).
      \]
      又
      \[
        \sum_{n=1}^{\infty}\frac{u^{n}}{n}=-\ln(1-u),\qquad |u|<1,
      \]
      令$u=\dfrac{x-3}{3}$，得
      \[
        S(x)=-\ln\left(1-\frac{x-3}{3}\right),\qquad x\in[0,6).
      \]
  \end{enumerate}

  \section{应用题(每小题 7 分，共 14 分)}
  \begin{enumerate}
    \setcounter{enumi}{16}
    \item \textbf{Solution}. 设所求点为$M(x,y,z)$，则
      \[
        \nabla f(M)=\{2x,2y,2z\},\qquad \boldsymbol{n}^{\circ}=\frac{1}{\sqrt2}\{0,-1,1\}.
      \]
      故方向导数为
      \[
        \frac{\partial f(M)}{\partial \boldsymbol{n}}=\sqrt2(-y+z).
      \]
      由 Lagrange 乘子法得$x=0,y=-z$，结合约束$x^{2}+2y^{2}+2z^{2}=1$，
      受检点为
      \[
        \left(0,\frac12,-\frac12\right),\qquad \left(0,-\frac12,\frac12\right).
      \]
      比较可知最大值点为
      \[
        \left(0,-\frac12,\frac12\right).
      \]

    \item \textbf{Solution}. 取底面$S_{1}:z=0,\ x^{2}+y^{2}\leq1$的下侧，使$S+S_{1}$封闭并取外侧法向.
      由高斯公式，
      \[
        \Phi=\iiint\limits_{V}(6x^{2}+6y^{2}+6z)\,\mathrm{d}v-\iint\limits_{S_{1}}3(z^{2}-1)\,\mathrm{d}x\,\mathrm{d}y.
      \]
      其中
      \[
        \iiint\limits_{V}(6x^{2}+6y^{2}+6z)\,\mathrm{d}v
        =6\int_{0}^{2\pi}\,\mathrm{d}\theta\int_{0}^{1}r\,\mathrm{d}r\int_{0}^{1-r^{2}}(r^{2}+z)\,\mathrm{d}z
        =2\pi,
      \]
      且
      \[
        \iint\limits_{S_{1}}3(z^{2}-1)\,\mathrm{d}x\,\mathrm{d}y=3\pi.
      \]
      因此
      \[
        \Phi=-\pi.
      \]
  \end{enumerate}

  \section{综合题(每小题 5 分，共 10 分)}
  \begin{enumerate}
    \setcounter{enumi}{18}
    \item \textbf{Proof}. 参见教材中莱布尼兹判别法的证明.

    \item \textbf{Solution}. 因$f(1,y)=0,f(x,1)=0$，所以
      \[
        f_{y}(1,y)=0,\qquad f_{x}(x,1)=0.
      \]
      由两次分部积分，
      \[
        \begin{aligned}
          I
          &=\iint\limits_{D}xyf_{xy}(x,y)\,\mathrm{d}\sigma
          =\int_{0}^{1}\,\mathrm{d}x\int_{0}^{1}xyf_{xy}(x,y)\,\mathrm{d}y \\
          &=\int_{0}^{1}\,\mathrm{d}x\left[xyf_{x}(x,y)\right]_{0}^{1}
          -\int_{0}^{1}\,\mathrm{d}x\int_{0}^{1}xf_{x}(x,y)\,\mathrm{d}y \\
          &=-\int_{0}^{1}\,\mathrm{d}y\left[xf(x,y)\right]_{0}^{1}
          +\int_{0}^{1}\,\mathrm{d}y\int_{0}^{1}f(x,y)\,\mathrm{d}x \\
          &=\iint\limits_{D}f(x,y)\,\mathrm{d}\sigma=a.
        \end{aligned}
      \]
  \end{enumerate}
\end{document}






